\chapter{Rigidity of entire graphical self-shrinkers}\label{sec:entire-graphical-shrinkers}
An entire graphical self-shrinker is more rigid than a general entire minimal graph. We develop the argument of \cite{Wang09}, Theorem 1.1 and Lemmas 2.1--2.3, in the normalization $H=\langle X,\nu\rangle/2$. For an upward graph this equation is
\[
 \left(\delta_{ij}-\frac{u_i u_j}{1+|Du|^2}\right)u_{ij}
       =\frac{x\cdot Du-u}{2}.
\]
The weighted argument, the height estimate and the area estimate are separate steps. The area estimate below uses a cylinder divergence calculation, so it does not require the intersection with a sphere to have a prescribed topology. We write $\rho=(4\pi)^{-n/2}e^{-|X|^2/4}$ and retain the Gaussian drift $\mathcal D$ of Definition \ref{def:gaussian-drift}.

\begin{proposition}[Positive angle gives a shifted Gaussian stability identity]\label{prop:shrinker-positive-angle-ground-state}
Let $X:\Sigma^n\to\mathbb R^{n+1}$ be a smooth two-sided self-shrinker, and suppose a constant vector $e$ has $w=\langle\nu,e\rangle>0$. For every smooth compactly supported function $\phi$ in the interior,
\[
 \int_\Sigma (|\nabla\phi|^2-|A|^2\phi^2)\rho
    =\int_\Sigma w^2|\nabla(\phi/w)|^2\rho\geq0.
\]
This is the quadratic form of $-(L-1/2)$, not of the fixed-center Gaussian Jacobi operator $-L$. No uniform positive lower bound for $w$ is assumed. Compare \cite{Wang09}, Lemma 2.1.
\end{proposition}
\begin{proof}
Proposition \ref{prop:shrinker-translation-dilation} gives $Lw=w/2$, hence $\mathcal Dw+|A|^2w=0$. Write $\phi=wz$, which is smooth and compactly supported because $w>0$ near its support. Expanding the square shows
\[
 |\nabla\phi|^2-w^2|\nabla z|^2
  =z^2|\nabla w|^2+2wz\langle\nabla w,\nabla z\rangle
  =\langle\nabla w,\nabla(wz^2)\rangle.
\]
Integrate using Lemma \ref{lem:gaussian-integration-parts}. The right-hand side becomes $-\int wz^2\mathcal Dw\rho=\int|A|^2\phi^2\rho$, proving the identity. Compactly supported Lipschitz tests follow by local smooth approximation. The distinction between $L$ and $L-1/2$ is essential: the identity does not contradict fixed-center instability.
\end{proof}

\begin{theorem}[Finite Gaussian area and a strict Gauss-map hemisphere force a plane]\label{thm:shrinker-positive-angle-rigidity}
Let a connected complete boundaryless smooth two-sided immersed self-shrinker have finite Gaussian area. If $\langle\nu,e\rangle>0$ everywhere for some constant vector $e$, then its image is a hyperplane through the origin. This extends the cutoff step of \cite{Wang09}, Lemma 2.1, to the stated immersed setting. Properness and a curvature bound are not additional hypotheses here.
\end{theorem}
\begin{proof}
Use intrinsic distance cutoffs $\eta_R$, equal to one on the radius-$R$ intrinsic ball about a fixed point, supported in the radius-$2R$ ball and with $|\nabla\eta_R|\leq2/R$. Completeness makes their supports compact by Hopf--Rinow. Proposition \ref{prop:shrinker-positive-angle-ground-state} gives
\[
 \int|A|^2\eta_R^2\rho\leq\frac4{R^2}\int\rho\longrightarrow0.
\]
Every fixed compact subset is eventually in the region where $\eta_R=1$. Nonnegativity and smoothness therefore force $A=0$ everywhere, without first assuming a global curvature integral. Lemma \ref{lem:complete-parallel-hypersurface} identifies the complete image as an affine plane. Its mean curvature is zero; the shrinker equation then makes its normal offset zero. No estimate for $1/w$ at infinity was used.
\end{proof}

\begin{lemma}[A compact flow avoids a smooth entire graph without growth bounds]\label{lem:compact-flow-entire-graph-avoidance}
Let $u:\mathbb R^n\times[0,T]\to\mathbb R$ be a smooth entire graphical mean curvature flow, and let $N_t$ be a smooth mean curvature flow of a compact boundaryless hypersurface on the same interval. If the two initial images are disjoint, their distance is positive and nondecreasing. No global slope, height or curvature bound for the entire graph is required.
\end{lemma}
\begin{proof}
The images of $N_t$ lie in one fixed compact set on a compact time interval. The distance to the graph is attained: a minimizing graph point cannot have arbitrarily large base coordinate, since horizontal distance alone would then tend to infinity. More uniformly, choose one graph base point and one parameter point of $N_t$ to obtain a finite upper bound for the minimizing distance over all times. This confines every minimizing graph base point to one fixed compact ball. Thus all minimizing pairs lie in a fixed compact parameter set on which both flows and their derivatives are bounded.

The closest-pair calculation of Proposition \ref{prop:mcf-avoidance} now applies on this compact set until any possible first contact. At a positive-distance minimizing pair, the displacement is normal to both hypersurfaces, so their tangent spaces coincide. Simultaneous variations along matching orthonormal tangent vectors give the nonnegative time derivative of squared distance. Tangential velocities have zero pairing with the displacement there. The minimum is locally Lipschitz in time, so its almost-everywhere derivative inequality makes it nondecreasing. Its initially attained positive value prevents first contact. The entire graph hypothesis supplies the needed compactness of minimizing pairs without a maximum principle at infinity. This is the avoidance mechanism used for the sphere barriers in \cite{Wang09}, Lemma 2.3.
\end{proof}

\begin{lemma}[Shrinker graphs have a linear height bound]\label{lem:shrinker-graph-linear-height}
If $u$ is a smooth entire graphical self-shrinker, let $r_*=2\sqrt n$ and $M_* =\max_{|x|\leq r_*}|u(x)|$. For every $R\geq1$,
\[
 \max_{|x|\leq R}|u(x)|
      \leq (\sqrt2 M_*+r_*+1)R.
\]
In particular this is a conclusion from the equation, not a growth assumption. The sphere-barrier method is \cite{Wang09}, Lemma 2.3; the displayed constant is a convenient coarse bound.
\end{lemma}
\begin{proof}
Set $\lambda(t)=\sqrt{R^2+1-t}$ and
$U(x,t)=\lambda(t)u(x/\lambda(t))$ for $0\leq t\leq R^2$.
Differentiation and the shrinker equation in the section introduction show that $U_t=a^{ij}(DU)U_{ij}$: both sides equal $(y\cdot Du(y)-u(y))/(2\lambda)$ with $y=x/\lambda$. Thus these are smooth entire graphical flow slices. Their initial height above the base ball of radius $r_*R$ is at most $\lambda(0)M_*$, because $r_*R/\lambda(0)<r_*$.

Place a sphere of radius $r_*R$ at height
$c=\lambda(0)M_*+r_*R+1$ above the origin in the base. Initially its entire ball lies strictly above the graph wherever their base projections overlap. Outside that projection no graph point is in the ball. The sphere evolves with radius
$r(t)=\sqrt{r_*^2R^2-2nt}$, which is at least $\sqrt{2n}R>R$ on the whole interval. Lemma \ref{lem:compact-flow-entire-graph-avoidance} keeps the sphere disjoint from the graph. For each fixed $|x|\leq R$, the graph height starts below the sphere's lower sheet and cannot cross it by continuity; this sheet exists at that base point for every time because $r(t)>R$. Hence $U(x,R^2)<c$. Since $U(x,R^2)=u(x)$ and $\lambda(0)\leq\sqrt2R$, this is the asserted upper bound. Apply the same construction to $-u$, which satisfies the same graph equation, for the lower bound. No estimate on $Du$ outside a compact ball enters the argument.
\end{proof}

\begin{proposition}[A cylinder area estimate for a graph with bounded mean curvature]\label{prop:graph-area-height-mean-curvature}
Suppose $u$ is smooth on a neighborhood of the closed base ball $\overline B_R^n$, $R>0$, with $|u|\leq M$ and $|H|\leq K$ there. Its graph area over $B_R^n$ satisfies
\[
 \operatorname{Area}(\operatorname{Graph}_{B_R^n}u)
 \leq\omega_nR^n+n\omega_nR^{n-1}(2M+R)
                    +K\omega_nR^n(2M+R),
\]
where $\omega_n$ is the volume of the Euclidean unit $n$-ball. The estimate does not assume a slope bound.
\end{proposition}
\begin{proof}
Extend the upward unit normal to the vertical cylinder by
$N(x,z)=(-Du(x),1)/\sqrt{1+|Du(x)|^2}$, independent of $z$. Its ambient divergence is $H(x)=-\operatorname{div}(Du/v)$. Apply the divergence theorem to
$E=\{(x,z):|x|<R,\ u(x)<z<M+R\}$.
This is a bounded piecewise smooth domain, with a positive vertical gap everywhere over the base. On its graph boundary the outward normal is $-N$, so its flux is minus the graph area. The top flux is at most the base volume because $|N|=1$. The lateral boundary has area at most $n\omega_nR^{n-1}(2M+R)$, and its flux is bounded by that area. Finally $|\int_E\operatorname{div}N|\leq K\operatorname{Vol}(E)\leq K\omega_nR^n(2M+R)$. Rearranging the flux identity gives the result. The codimension-two corners do not contribute surface flux; equivalently one may smooth them and pass to the limit. This is a cylinder version of the bounded-divergence comparison underlying \cite{Wang09}, Lemma 2.2.
\end{proof}

\begin{corollary}[Polynomial area and finite Gaussian area are automatic for shrinker graphs]\label{cor:shrinker-graph-polynomial-area}
The graph $\Sigma$ of every smooth entire self-shrinker satisfies
\[
 \mathcal H^n(\Sigma\cap B_R^{n+1})\leq C R^{n+2}\quad(R\geq1)
\]
for a finite constant depending only on $n$ and the height bound in Lemma \ref{lem:shrinker-graph-linear-height}. In particular it has finite Gaussian area. The exponent here is a sufficient polynomial bound; it is not asserted to be sharp.
\end{corollary}
\begin{proof}
On the graph above $B_R^n$, the shrinker equation gives
$|H|=|\langle(x,u),\nu\rangle|/2\leq\sqrt{R^2+M_R^2}/2$,
where $M_R=\max_{B_R^n}|u|\leq C_1R$. Insert $M=C_1R$ and $K=\sqrt{1+C_1^2}\,R/2$ in Proposition \ref{prop:graph-area-height-mean-curvature}. The first two terms are of order $R^n$ and the last of order $R^{n+2}$. The part of the graph in an ambient radius-$R$ ball is contained in the graph over the base radius-$R$ ball, so the same upper bound applies. Lemma \ref{lem:gaussian-polynomial-moments} now gives finite Gaussian area. This proves the growth input before the weighted cutoff argument, rather than assuming it in the final Bernstein theorem.
\end{proof}

\begin{theorem}[Every smooth entire graphical self-shrinker is a linear hyperplane]\label{thm:entire-shrinker-graph-bernstein}
If $u:\mathbb R^n\to\mathbb R$, $n\geq1$, is smooth and its graph satisfies $H=\langle X,\nu\rangle/2$, then $u(x)=p\cdot x$ for a constant vector $p\in\mathbb R^n$. No height, slope, curvature or volume-growth bound is assumed. This is \cite{Wang09}, Theorem 1.1, arXiv v1, with the growth and stability steps expanded below and above.
\end{theorem}
\begin{proof}
An entire smooth graph is connected, boundaryless and two-sided with strictly positive vertical angle. It is complete: a Cauchy sequence for its induced metric projects to a Euclidean Cauchy sequence because the metric dominates the base metric. The base limit is finite; smoothness of $u$ makes the induced and base metrics uniformly comparable on a small neighborhood there, so the original sequence converges on the graph. Corollary \ref{cor:shrinker-graph-polynomial-area} supplies finite Gaussian area. Theorem \ref{thm:shrinker-positive-angle-rigidity} therefore makes the image a hyperplane through the origin. Because it is an entire graph, this plane is the graph of a linear function. Alternatively $A=0$ gives $D^2u=0$, so $u=p\cdot x+b$, and the graph shrinker equation gives $b=0$. This proof does not use the dimension-dependent Bernstein theorem for minimal graphs or a global parabolic gradient estimate.
\end{proof}

\begin{example}[A spherical cap has positive angle but is incomplete]\label{ex:shrinker-hemisphere-incomplete}
The graph
\[
 u(x)=\sqrt{2n-|x|^2},\qquad |x|<\sqrt{2n},
\]
is the upper hemisphere of the round self-shrinker. Its upward angle is $w=u/\sqrt{2n}>0$, its squared curvature is $1/2$, and its Gaussian area is finite. It is incomplete and its graph domain is bounded, so it satisfies neither the complete-image hypothesis of Theorem \ref{thm:shrinker-positive-angle-rigidity} nor the entire-domain hypothesis of Theorem \ref{thm:entire-shrinker-graph-bernstein}.
\end{example}
\begin{proof}
The round sphere of radius $\sqrt{2n}$ is a self-shrinker by Proposition \ref{prop:shrinking-cylinders}. On its upper hemisphere the outward normal is $X/\sqrt{2n}$ and points upward, giving the angle. Its $n$ principal curvatures equal $1/\sqrt{2n}$, so their squared sum is $1/2$. Its area and Gaussian weight are finite. A meridian reaches the omitted equator in finite intrinsic length, proving incompleteness. In particular one cannot insert intrinsic exhaustion cutoffs with compact support and gradient tending uniformly to zero as in the complete proof.
\end{proof}
